% ============================================================================
%  sections/3_requirement_analysis.tex -- Chapter 3: REQUIREMENT ANALYSIS
% ============================================================================
\chapter{REQUIREMENT ANALYSIS}
\section{Hardware Requirements}
The hardware components must support real-time environmental monitoring, local edge-based acoustic processing, and low-power long-range communication. The specific hardware requirements include:

\begin{itemize}
    \item \textbf{Microcontroller (ESP32-S3):} 
    Must feature a dual-core Xtensa LX7 processor running up to 240 MHz with vector instructions (to accelerate TinyML neural network calculations). It requires at least 512 KB of on-chip SRAM and support for external PSRAM to handle audio buffering and spectrogram conversion.
    
    \item \textbf{LoRa Transceiver (SX1278):} 
    Must operate in the 433 MHz ISM band with a high link budget (up to 170 dB) and high sensitivity to ensure reliable, multi-hop data transmission through dense forest canopies and rugged terrain.
    
    \item \textbf{Acoustic Sensor (MEMS Microphone):} 
    An omnidirectional digital MEMS microphone (such as the INMP441) using an I2S interface to stream high-fidelity, low-noise audio directly to the ESP32-S3 for acoustic analysis.
    
    \item \textbf{Environmental Sensors:} 
    An integrated array containing a high-accuracy temperature and humidity sensor (e.g., DHT22 or SHT3x) and a gas sensor (e.g., MQ series) to gather ambient environmental data for fire risk estimation.
    
    \item \textbf{Power Management System:} 
    A low-dropout (LDO) voltage regulator, a Li-ion/LiPo battery charger circuit, and a solar panel interface to sustain continuous, autonomous outdoor operations.
\end{itemize}

\section{Software Requirements}
The software stack must support local signal processing, machine learning inference, network routing simulation, and data visualization. The software requirements are classified as follows:

\begin{itemize}
    \item \textbf{Embedded Development Environment:} 
    \begin{itemize}
        \item \textbf{ESP-IDF / Arduino IDE:} Used for writing, compiling, and flashing the firmware onto the ESP32-S3.
        \item \textbf{TinyML Framework (TensorFlow Lite for Microcontrollers):} Required to load, optimize, and run the quantized acoustic classification model directly on the microcontroller.
    \end{itemize}
    
    \item \textbf{Digital Signal Processing (DSP) Libraries:}
    Libraries used to transform raw audio streams into Mel-spectrograms for
    model inference.
    
    \item \textbf{Simulation Tools:} 
    \begin{itemize}
        \item \textbf{OMNeT++:} Used to model the physical wireless network.
        \item \textbf{FLORA Framework:} An OMNeT++ based simulator utilized to simulate the LoRa physical layer, evaluate packet delivery ratios, and test the LDSE multi-hop routing protocol.
        \item \textbf{ns-3 Network Simulator:} Utilized to simulate the LoRa network, evaluate packet delivery ratio, latency, energy consumption, and network scalability, and validate the performance of the proposed LDSE multi-hop routing protocol.
    \end{itemize}
    
    \item \textbf{Application and Visualization Software:} 
    \begin{itemize}
        \item \textbf{Web Dashboard Framework:} A web-based application (built using modern web technologies) to display real-time sensor readings, historical trends, acoustic classification events, and node health status.
    \end{itemize}
\end{itemize}
\section{Feasibility Study}
\subsection{Technical Feasibility:}The ESP32-S3 microcontroller provides sufficient computational capability for TinyML inference at the edge, with its Xtensa LX7
dual-core processor and a vector instruction set enabling efficient neural network
computation within the available 512~KB SRAM. The SX1278 LoRa transceiver provides
long-range communication with a link budget exceeding 170~dB, well suited for the
dense vegetation and terrain variation found in Nepalese community forests. The
Edge Impulse platform provides a mature and well-documented pipeline for training,
quantizing, and deploying CNN models to the ESP32-S3 target, and TensorFlow Lite
Micro has been deployed successfully on similar hardware in prior academic work.
All hardware components selected for this project are commercially available and
have documented community support.

\subsection{Economic Feasibility:} The estimated total prototype budget of
approximately NPR 12,400 positions this system as a highly cost-effective
alternative to commercial forest monitoring solutions, which typically require
expensive satellite or cellular infrastructure.

The use of commodity hardware components and open-source software stacks
including Edge Impulse, TensorFlow Lite Micro, and the Arduino framework
eliminates licensing costs. The modular and scalable architecture means that
additional forest nodes can be added to expand coverage at a marginal per-node
cost well below the gateway cost, making the system economically viable for
community forest committees operating under budget constraints.

\subsection{Operational Feasibility:} The system is designed to operate autonomously
with deep-sleep scheduling and event-driven wake-up, requiring minimal human
intervention after initial deployment and configuration. Solar-assisted battery
operation can extends node lifetime significantly in remote forest environments
where regular maintenance visits are impractical. The real-time web dashboard
provides an intuitive interface for forest rangers to monitor alerts, review event
histories, and identify geographic patterns of illegal activity without requiring
technical expertise. Fail-safe design including confidence thresholding for alert
transmission reduces false alarm rates, maintaining the trust of operators in the
system's outputs.

\subsection{Time Feasibility:} The proposed two month project schedule is realistic
given the availability of the FSC22 dataset, the Edge Impulse platform's rapid
model development capability, the team's familiarity with embedded systems and IoT
development, and the modular architecture of the system design. Each phase builds
incrementally on the previous one, and the independence of the acoustic model
training, fire-risk model development, LDSE protocol implementation, and dashboard
development workstreams allows parallel progress across team members.

\clearpage